% ============================================================
% Capítulo 13 — Apéndices
% ============================================================
\chapter{Apéndices}
\label{cap:apendices}

\section{Glosario}

\begin{description}[leftmargin=2.5cm,style=nextline]
  \item[\textbf{BERT}] \emph{Bidirectional Encoder Representations
    from Transformers}. Modelo de lenguaje preentrenado de Google
    (2018). En este proyecto se usa sólo como tokenizer.

  \item[\textbf{BiGRU}] GRU bidireccional. Variante de RNN que
    procesa la secuencia en ambas direcciones y concatena los
    estados ocultos.

  \item[\textbf{CBIR}] \emph{Content-Based Image Retrieval}.
    Recuperación de imágenes por contenido visual en lugar de
    metadatos.

  \item[\textbf{CMC}] \emph{Cumulative Matching Characteristic}.
    Curva que mide la fracción de queries cuyo positivo aparece
    en el top-k. \texttt{cmc[k]} = Top-(k+1) accuracy.

  \item[\textbf{CUHK-PEDES}] \emph{CUHK Person Description Dataset}.
    Benchmark estándar de \reid{} por texto, 40\,206 imágenes con
    descripciones.

  \item[\textbf{DSC}] \emph{Depthwise Separable Convolution}.
    Factorización de una convolución estándar en depthwise
    (por canal) + pointwise ($1 \times 1$). Reduce parámetros y
    FLOPs.

  \item[\textbf{EfficientNet-B0}] Familia de CNNs con compound
    scaling. B0 es la variante más pequeña (\textasciitilde 5,3M
    parámetros), preentrenada en ImageNet.

  \item[\textbf{FPN}] \emph{Feature Pyramid Network}. Estructura
    que combina features de múltiples escalas para detección.

  \item[\textbf{GRU}] \emph{Gated Recurrent Unit}. RNN con dos
    compuertas (reset, update), más simple que LSTM.

  \item[\textbf{mAP}] \emph{mean Average Precision}. Promedio de
    AP sobre todas las queries. Mide calidad del ranking
    completo.

  \item[\textbf{NMS}] \emph{Non-Maximum Suppression}. Técnica para
    eliminar detecciones redundantes.

  \item[\textbf{Person Re-ID}] Tarea de reconocer a la misma
    persona en diferentes cámaras o momentos.

  \item[\textbf{Triplet loss}] Pérdida que enforce que un anchor
    esté más cerca de un positivo que de un negativo, con un
    margen dado.

  \item[\textbf{U-TextReIDNet}] \emph{Unified Text ReID Network}.
    Modelo del paper que une detección humana con \reid{} por
    texto.

  \item[\textbf{Top-k}] Fracción de queries cuyo positivo aparece
    entre los primeros $k$ resultados.
\end{description}

\section{Tabla maestra de parámetros del modelo}

\begin{longtable}{p{4cm}p{3.5cm}p{6cm}}
\toprule
\textbf{Capa} & \textbf{Forma de salida} & \textbf{Parámetros} \\
\midrule
\endhead
\texttt{Embedding(29610, 512)}
  & \dims{(B, 100, 512)}
  & 15,16 M \\
\texttt{Dropout(0.30)}
  & \dims{(B, 100, 512)}
  & 0 \\
\texttt{BiGRU(512, 1024, bidir)}
  & \dims{(B, 100, 2048)}
  & 4,72 M \\
Suma fwd+bwd / 2
  & \dims{(B, 100, 1024)}
  & 0 \\
Permute + unsqueeze
  & \dims{(B, 1024, 100, 1)}
  & 0 \\
\midrule
\texttt{EfficientNet-B0}
  & \dims{(B, 1280, 16, 16)}
  & 5,29 M \\
\texttt{DSC(1280, 1024)}
  & \dims{(B, 1024, 16, 16)}
  & \textasciitilde 14 K \\
\texttt{AdaptiveMaxPool2d(1, 1)}
  & \dims{(B, 1024, 1, 1)}
  & 0 \\
\texttt{DSC(1024, 1024)}
  & \dims{(B, 1024, 1, 1)}
  & \textasciitilde 10 K \\
\midrule
Classifier(1024, num\_classes)
  & \dims{(B, num\_classes)}
  & \textasciitilde 11 K por clase \\
\midrule
\textbf{Total (sin clasificador)}
  & ---
  & \textbf{\textasciitilde 25,2 M} \\
\textbf{Total (con clasificador)}
  & ---
  & \textbf{\textasciitilde 25,2 M + $1024 \cdot N_c$} \\
\bottomrule
\caption{Recuento completo de parámetros. $N_c$ = número de
identidades (clases) en el dataset de entrenamiento. Para
CUHK-PEDES, $N_c = 11\,003$.}
\label{tab:master-params}
\end{longtable}

\section{Tabla maestra de hiperparámetros de entrenamiento}

\begin{longtable}{p{4cm}p{4cm}p{6cm}}
\toprule
\textbf{Hiperparámetro} & \textbf{Valor} & \textbf{Ubicación} \\
\midrule
\endhead
Épocas
  & 60
  & \fline{config.py}{71} \\
Batch size (default)
  & 8
  & \fline{config.py}{123,129} \\
Batch size (Jetson)
  & 4
  & \fline{config.py}{126} \\
Learning rate inicial
  & 1e-3
  & \fline{config.py}{75} \\
Weight decay
  & (decoupled, AdamW)
  & \fline{train.py}{98} \\
$\beta_1$ AdamW
  & 0.90
  & \fline{config.py}{72} \\
$\beta_2$ AdamW
  & 0.999
  & \fline{config.py}{73} \\
LR decay epochs
  & [20, 40]
  & \fline{config.py}{74} \\
LR decay factor
  & 0.1 (default PyTorch)
  & (implícito) \\
Mixed precision (GPU)
  & FP16
  & \fline{train.py}{130} \\
Mixed precision (CPU)
  & BF16
  & \fline{train.py}{130} \\
Ranking loss weight $\alpha$
  & 1.0
  & \fline{config.py}{87} \\
Identity loss weight $\beta$
  & 1.0
  & \fline{config.py}{88} \\
Triplet margin
  & 0.5
  & \fline{config.py}{86} \\
Seed
  & 3407
  & \fline{config.py}{79} \\
\bottomrule
\caption{Hiperparámetros de entrenamiento.}
\label{tab:master-hyperparams}
\end{longtable}

\section{Resultados reportados en el paper}

\begin{longtable}{p{3.5cm}p{2.5cm}p{2.5cm}p{2.5cm}}
\toprule
\textbf{Dataset} & \textbf{Top-1} & \textbf{Top-5} & \textbf{Top-10} \\
\midrule
\endhead
CUHK-PEDES & 54,02\% & 75,51\% & 85,28\% \\
RSTPReid   & 38,45\% & ---    & ---    \\
\bottomrule
\caption{Resultados del paper original.}
\label{tab:paper-results}
\end{longtable}

\noindent\textbf{U-TextReIDNet}: 38,77 M parámetros totales.

\section{Mapa de archivos}

\begin{verbatim}
textreid-train/
|-- README.md
|-- config.py                     # Configuración global
|-- train.py                      # Script de entrenamiento
|-- evaluate.py                   # Script de evaluación
|-- requirements.txt
|
|-- model/
|   |-- textreidnet.py            # Modelo principal
|   |-- visual_network.py         # EfficientNet-B0 wrapper
|   |-- language_network.py       # BERT tokenizer + BiGRU
|   |-- efficientnet_backbone.py  # Variante multi-stage (no usada)
|   |-- model_utils.py            # DepthwiseSeparableConv, etc.
|
|-- datasets/
|   |-- bases.py                  # ImageTextDataset, etc.
|   |-- cuhkpedes.py              # Loader dataset original
|   |-- cuhkpedes_hf.py           # Loader dataset HuggingFace
|   |-- cuhkpedes_dataloader.py   # Constructor de DataLoaders
|   |-- bert_tokenizer.py         # BERT tokenizer
|   |-- simple_tokenizer.py       # CLIP-style BPE tokenizer
|   |-- tiktoken_tokenizer.py     # OpenAI tiktoken
|
|-- evaluation/
|   |-- ranking_loss.py           # Triplet semi-hard
|   |-- identity_loss.py          # Cross-entropy simétrica
|   |-- focal_loss.py             # Sigmoid focal (no usada)
|   |-- evaluations.py            # Top-k, mAP
|
|-- utils/
|   |-- iotools.py                # read_image, read_json
|   |-- miscellaneous_utils.py    # set_seed, pad_tokens, etc.
|
|-- scripts/
|   `-- download_hf_dataset.py    # Descarga y convierte HF
|
|-- data/
|   `-- CUHK-PEDES-HF/            # Dataset descargado
|       |-- raw.parquet
|       |-- reid_raw.json
|       `-- imgs/all/*.jpg
|
|-- docs/
|   `-- paper.md                  # Paper traducido al español
|
|-- notebooks/
|   `-- 01_lab_validate_dataset.ipynb
|
`-- book/                         # (este libro)
    |-- main.tex
    |-- chapters/*.tex
    `-- figures/*.tikz
\end{verbatim}

\section{Cheatsheet de Bash}

\begin{minted}[language=bash, basicstyle=\ttfamily\small, frame=single]{python}
# Instalar dependencias
pip install -r requirements.txt --extra-index-url https://download.pytorch.org/whl/cu117

# Descargar dataset
python scripts/download_hf_dataset.py

# Entrenar
python train.py --dataset_source huggingface --epochs 60

# Evaluar
python evaluate.py \
    --checkpoint data/checkpoints/TextReIDNet_latest.pth.tar \
    --dataset_source huggingface \
    --split test

# Inspeccionar configuración
python -c "from config import sys_configuration; \
    cfg = sys_configuration(); print(cfg)"
\end{minted}

\section{Referencias}

\begin{itemize}
  \item \textbf{Agyeman, R. \& Rinner, B.} (2024).
    \emph{Decentralized Text-Based Person Re-Identification in
    Multicamera Networks}. IEEE Access.
    DOI: 10.1109/ACCESS.2024.3501382.
    Traducción al español en \file{docs/paper.md}.

  \item \textbf{Tan, M. \& Le, Q. V.} (2019).
    \emph{EfficientNet: Rethinking Model Scaling for
    Convolutional Neural Networks}. ICML.
    arXiv:1905.11946.

  \item \textbf{Devlin, J., Chang, M.-W., Lee, K. \& Toutanova,
    K.} (2018). \emph{BERT: Pre-training of Deep Bidirectional
    Transformers for Language Understanding}. NAACL.
    arXiv:1810.04805.

  \item \textbf{Cho, K. et al.} (2014). \emph{Learning Phrase
    Representations using RNN Encoder-Decoder for Statistical
    Machine Translation}. EMNLP.

  \item \textbf{Howard, A. G. et al.} (2017).
    \emph{MobileNets: Efficient Convolutional Neural Networks
    for Mobile Vision Applications}. arXiv:1704.04861.

  \item \textbf{Li, S., Xiao, T., Li, H., Zhou, B., Yue, D. \&
    Wang, X.} (2017). \emph{Person Search with Natural Language
    Description}. ICCV.

  \item \textbf{Loshchilov, I. \& Hutter, F.} (2019).
    \emph{Decoupled Weight Decay Regularization}. ICLR.
    (AdamW).

  \item \textbf{Lin, T.-Y., Goyal, P., Girshick, R., He, K. \&
    Dollár, P.} (2017). \emph{Focal Loss for Dense Object
    Detection}. ICCV. (RetinaNet).
\end{itemize}
